\begin{abstract}
\label{sec:abstract}
语言模型拒答行为的权重干预需要同时回答两个问题：局部表示如何改变，改变后能否满足预设的行为与分布约束。本文以 Heretic 方向消融为起点，在注明上游任意秩消融来源的基础上，研究本地校准式实现及其后续适配。单位置版本冻结基座输入输出，通过秩上界为 128 的增量、尺度归一化的平滑邻居距离、保持与间隔推远目标及因子规范化优化局部模块；轨迹版本进一步实现逐步对齐采集、部署增益和隔离审计协议。本文离线重放一次 Qwen3.8-27B 归档搜索的全部 120 个候选，全部状态为完成，但无候选通过验收门槛。验证集最低拒答关键词命中率为 54/100，对应首 token KL 为 0.089975 nats；该候选并非独立测试结果或已验收模型。所有候选满足 KL 数值门槛，却均未满足关键词门槛，独立审计及重载分数无记录。结果说明，这次单位置搜索产生了局部行为变化，但未达到既定干预目标；低首 token 漂移不能替代完整能力评价。轨迹版本仅报告已实现机制，本归档没有其效果数据。本文据此给出参数化、数值实现与证据边界一致的方法分析，为后续同协议对照提供可复查基础。

\noindent\textbf{关键词：}机制可解释性；拒答行为；任意秩消融；低秩适配；实验复核
\end{abstract}
